\section{A second label tier: binary critical-residue maps}

The $\ddG$ antigens above are the gold standard --- a graded, energetic importance per
residue --- but experimental alanine scans with full $\ddG$ tables are scarce. A much larger
body of data exists at a coarser but still \emph{energetic} resolution: \textbf{saturating
shotgun alanine scans} read out by flow cytometry (Integral Molecular's platform; Davidson \&
Doranz), which classify each scanned residue as \emph{critical} or \emph{non-critical} for a
given antibody. A residue is critical when its mutation measurably abolishes binding, so even
this binary call is grounded in binding energy, not mere proximity --- it is the same
quantity as $\ddG$, thresholded. We therefore treat these as a second label tier:
\textbf{Tier~1}, graded $\ddG$ (\S\ref{sec:simulated-antigens}); \textbf{Tier~2}, binary
critical/non-critical.

\paragraph{Source and the proximity filter.} We mine these from IEDB
(\texttt{query-api.iedb.org}), keeping only conformational (discontinuous) epitopes defined by
\emph{binding} assays and explicitly discarding every structure/contact-derived record --- the
proximity labels this project argues against (recall that $>99\%$ of structure-linked IEDB
epitopes are crystal/cryo-EM contact, \S\ref{sec:motivation}). Within the binding-derived set
we take the flow-cytometry (shotgun-mutagenesis) subset, whose curated residue list \emph{is}
the set of critical residues. This yields $3{,}213$ antibody--antigen epitope maps over
$186$ antigens. Table~\ref{tab:binary-survey} lists the simulatable, not-yet-included
candidates; note the depth of antibody sampling per antigen (e.g.\ dengue envelope: $781$
mapped antibodies across $30$ studies), which lets a Tier-2 antigen carry a
\emph{consensus} epitope (residues critical across many antibodies) as well as
antibody-specific maps.

% AUTO-GENERATED by analysis/make_binary_label_table.py -- do not edit by hand.
\begin{table}[htbp]\centering\small
\caption{\textbf{Tier-2 (binary) label candidates} from IEDB shotgun-mutagenesis alanine scans (flow-cytometry critical-residue mapping; Integral Molecular / Doranz). Labels are per-residue \emph{critical vs non-critical}, energetically grounded (a critical residue is one whose mutation measurably weakens antibody binding) but coarser than $\ddG$. ``Antibodies'' = mapped antibody footprints (one antibody $\to$ one conformational epitope); ``Studies'' = distinct source papers; ``Crit.\ res.'' = union of critical residues across all antibodies. Structures provisional ([v] = numbering to be verified before use).}
\label{tab:binary-survey}
\begin{tabular}{llrrrl}
\toprule
Antigen & Organism & Antibodies & Studies & Crit.\ res. & Structural handle \\
\midrule
Dengue E (envelope) & Dengue virus & 781 & 30 & 632 & sE ectodomain, e.g.\ 1OKE [v] \\
Chikungunya E1/E2 & Alphavirus chikungunya & 128 & 6 & 98 & ectodomain, e.g.\ 3N42 [v] \\
RVFV Gn & Rift Valley fever virus & 63 & 8 & 210 & Gn ectodomain, e.g.\ 6F9F [v] \\
Influenza HA & Influenza A virus & 81 & 22 & 296 & HA ectodomain (trimeric) [v] \\
GRO-$\alpha$ / CXCL1 & Homo sapiens (human) & 30 & 1 & 23 & chemokine, e.g.\ 1MGS [v] \\
CD4 & Homo sapiens (human) & 1 & 1 & 18 & D1--D2, e.g.\ 1CDH [v] \\
CD40L & Homo sapiens (human) & 1 & 1 & 13 & TNF domain, e.g.\ 1ALY [v] \\
IL-2R$\alpha$ & Homo sapiens (human) & 2 & 1 & 13 & ectodomain, e.g.\ 1Z92 [v] \\
IL-7R$\alpha$ & Homo sapiens (human) & 2 & 1 & 32 & ectodomain, e.g.\ 3DI2 [v] \\
HLA class I & Homo sapiens (human) & 15 & 5 & 38 & peptide-MHC, many [v] \\
\bottomrule
\end{tabular}
\end{table}


\paragraph{Representation.} Each Tier-2 antigen is rendered like the $\ddG$ antigens but with a
\emph{categorical} colouring --- red $=$ critical, blue $=$ scanned but not critical, grey $=$
not scanned --- so the two tiers read consistently while remaining visually distinct
(Fig.~\ref{fig:binary-flagship}). The shotgun residues are reported in UniProt numbering and
carry their wild-type identity, so they are mapped onto a simulation structure and
wild-type-verified before use, exactly as for the $\ddG$ antigens. Tier-2 antigens enter the
benchmark as a separate label stream (binary), never merged into the $\ddG$ tables; in the
leakage-controlled evaluation they are held out by antigen like everything else.

\begin{figure}[H]\centering\figorbox{binary_flagship_ddg.png}{0.5}
\caption{\textit{Placeholder.} Flagship Tier-2 critical-residue map (red $=$ critical, blue $=$
scanned/non-critical, grey $=$ not scanned), pending numbering verification of the chosen
antigen against its simulation structure.}\label{fig:binary-flagship}
\end{figure}
